\scanpage{381}
\mspage{297}
% archive sheet number: 380

Donc le composé
\diagram{\begin{tikzcd}[ampersand replacement=\&]
J \arrow[r] \& J^B \arrow[r,"{\mathrm{cof.}}"] \& I^B
\end{tikzcd}}
est cofinal. Mais comme \deleted{il s'agit}\inserted{le deuxième}
\deleted{de} foncteur \inserted{$J^B\unsure{\text{---}|\longrightarrow}I^B$ est} pl. fid., \deleted{la catégorie}
il s'ensuit aussitôt que $J\longrightarrow J^B$
est cofinal (car $\unsure{\xi|\zeta}\backslash J$ pour $\xi\in\operatorname{Ob}J^B$
est le $=$ \inserted{à iso près}, qu'on considère $\xi$ comme
objet de $J^B$, ou de $I^B$).

Ainsi dans l'énoncé de la condition
$d)_{\unsure{r|\alpha}}$ du lemme, il est inutile de
spécifier que $J_r$ est g\textsuperscript{r}de devant $\pi$,
et de $=$ dans la conclusion sur $J$.

Comme $J\hookrightarrow I$ est pl. fidèle,
pour voir que ce foncteur est cofinal,
il suffit de voir qu'il satisfait F1,
i.e. que pour tout $i\in\operatorname{Ob} I$,
existe $j\in J$ avec $i\longrightarrow j$. On
va construire $j$ comme limite d'un
\deleted{suite \unsure{ord.}}\inserted{syst. inductif.} ordin\unsure{al|ale} d'objets \inserted{$i_\alpha$} de $I$,
indexée par des ordinaux $\alpha$ tels que
$\operatorname{card}\alpha<\pi\deleted{\unsure{{}_0|{}'_0}}$.
\deleted{On aura \unsure{donc} \inserted{\unsure{Comme} $\pi'_0<\pi$},} Désignant
par $\alpha_\pi$ le plus petit ordinal tel que
